\subsection{Polynomial vacuum}\label{c01-c1-polynomial-vacuum}

Polynomial exchange and fusion constructions occur in the dense and dilute $O(1)$ models~\cite{DiFrancescoZinnJustin2005,GarbaliNienhuis2017Ground,GarbaliNienhuis2017Sum}. We construct the present normalized vacuum directly by polynomial interpolation and determine its empty component. Its Pfaffian formula supplies the denominator of the polygon observable.

\begin{proposition}[Polynomial vacuum]\label{prop:pv:vacuum}
Use twist \(\zeta=t^2=i\). Let \(\rho X=(x_N,x_1,\ldots,x_{N-1})\), and \(\Pi_\kappa\) the corresponding rotation of spin slots with factor \(\kappa^{h_N}\) from the last (now first) spin. We construct a charge-zero vector polynomial \(\Psi_N(X)\) with
\[
 F_{i,i+1}(x_i/x_{i+1})\Psi(X)=\Psi(s_i X),\qquad
 \Pi_\zeta\Psi(X)=\Psi(\rho X)
\]
(\(s_i\) adjacent transposition), homogeneous of degree \(\binom N2\),
with degree at most \(N-1\) in each variable. For a spin component
\(\Psi_{N,\boldsymbol h}\), the refined bounds are
\[
 h_j=+1\ \Longrightarrow\ \deg_{x_j}\Psi_{N,\boldsymbol h}\le N-2,
 \qquad
 h_j=-1\ \Longrightarrow\ x_j\mid\Psi_{N,\boldsymbol h}.
\]
It has the adjacent-pair reduction
\[
 \Psi_N(\ldots,a,Ta,\ldots)=s_C(a;Y)\ \iota_C \Psi_{N-2}(Y),
 \qquad
 s_C(a;Y)=(1+T)a\prod_{x\in Y}(x-t^5a)(x-t^6a),
\]
where \(Y\) is the remaining list and \(\iota_C\) inserts a \(C\) tensor in the indicated positions; and at last parameter zero its projection onto last spin zero is \((\prod_{j<N}x_j)\Psi_{N-1}\).

\end{proposition}

\begin{proof}
We induct on \(N\). The construction has five steps: prescribe cup
values and interpolate; prove the spectator degree bounds; propagate
all cup reductions; prove the last exchange; and prove cyclicity.
The refined bounds are the reason the last two discrepancies vanish.

The initial vectors are \(\Psi_0=1\) and \(\Psi_{1,0}=1\), with
the other size-one components zero. At sizes two and three, the empty
components are
\[
 S_2(a,b)=a+b,\qquad
 S_3(a,b,c)=\sum_{i\ne j}x_i^2x_j+2\sqrt2\,abc,
 \quad (x_1,x_2,x_3)=(a,b,c).
\]
At size three, components with first pair \((h,-h)\) and last spin
zero are
\[
\begin{gathered}
 f_h(a,b;c)=2\sin\lambda\,t^h L_h(a,b)P(a,b;c),\\
 L_+(a,b)=b,\quad L_-(a,b)=a,\qquad
 P(a,b;c)=(c+ta)(c+b/t).
\end{gathered}
\]
At size two omit \(P\). Cyclicity specifies the remaining components.
We verify these initial cases at the end of the proof.

\paragraph{Cup values and interpolation.}
Assume the proposition at lower sizes and let \(N\ge4\). Set
\(Y_i=(x_1,\ldots,\widehat{x_i},\ldots,x_{N-1})\). To prescribe
the value at \(x_N=Tx_i\), insert \(C\) in the parameter list
\((x_1,\ldots,x_i,Tx_i,x_{i+1},\ldots,x_{N-1})\), then carry
its second slot to the end. With rightmost factor acting first, this is
\[
 V_i=s_C(x_i;Y_i)
 F_{N-1,N}(Tx_i/x_{N-1})\cdots
 F_{i+1,i+2}(Tx_i/x_{i+1})\,
 \iota_{C,i}\Psi_{N-2}(Y_i).
\]
The product is empty when \(i=N-1\). At the crossing with \(x_d\),
the two denominator factors are proportional to
\((x_d-t^5x_i)(x_d-t^6x_i)\), already present in \(s_C\).
Thus each \(V_i\) is polynomial.

These prescriptions are covariant under exchanges among the first
\(N-1\) sites, carrying the marked value with its site. An exchange
away from the mark uses lower-size exchange and Yang--Baxter. At the
mark, the local identity is the cup slide
\[
 F_{12}(a/b)F_{23}(Ta/b)(C\otimes1)=1\otimes C:
\]
it moves the cup based at \(a\) past its neighbor \(b\), without
changing \(s_C(a;Y)\). The inverse slide handles the opposite order;
Yang--Baxter moves a distant exchange into this local configuration.

For last spin plus, interpolate the \(N-1\) values with degree
\(N-2\). For last spin minus, interpolate the value divided by
\(x_N\) with that degree. For last spin zero, add the prescribed
zero node and use degree \(N-1\). More explicitly, put
\[
 \ell_i=\prod_{\substack{d<N\\d\ne i}}
 \frac{x_N-Tx_d}{T(x_i-x_d)}.
\]
The multiplier of \(V_i\) is \(\ell_i\) for last spin plus and
\(x_N\ell_i/(Tx_i)\) otherwise. In last spin zero add
\(\prod_{d<N}(x_d-x_N/T)\Psi_{N-1}\).

The apparent poles at \(x_i=x_d\) cancel: coincident nodes have
identical prescriptions by exchange covariance. Indeed exchanging
equal parameters through intervening slots is the identity by
unitarity, initially at generic other parameters. The only other
possible poles, the simple poles at \(x_i=0\), cancel because
\(x_i\mid V_i\). The result is polynomial, with total degree
\(\binom N2\) by homogeneity. Interpolation also proves all exchanges
among the first \(N-1\) sites.

\paragraph{Spectator bounds.}
Fix \(j<N\). For \(V_i\) with \(i\ne j\), the lower-size bound
and the quadratic factor in \(s_C\) give degree at most \(N-1\)
in \(x_j\). If slot \(j\) is crossed, then as \(x_j\to\infty\)
the crossing argument tends to zero. Its leading output spin at that
site is no larger than the input spin, so an output plus spin retains
the lower-size degree loss. At \(x_j=0\) the opposite triangularity
preserves the required minus-spin zero. The Lagrange multipliers are
bounded at both ends. The extra zero-node term has the same bounds
directly from size \(N-1\).

The term \(V_j\) has larger raw degree because \(x_j\) is the cup
parameter. Its orders and the two possible multipliers are
\[
\begin{array}{c|cc}
 &x_j\to\infty&x_j\to0\\ \hline
 V_j&O(x_j^{2N-3})&O(x_j)\\
 \ell_j&O(x_j^{-(N-2)})&O(1)\\
 x_N\ell_j/(Tx_j)&O(x_j^{-(N-1)})&O(x_j^{-1})
\end{array}
\]
These give all unrefined bounds. There are only two extra cases to
check. At infinity, if both site \(j\) and the final site have spin
plus, the carried cup spin starts minus and cannot increase under
the leading \(F(\infty)\) transport. The leading coefficient is
therefore zero, supplying one more degree of decay. At zero, if
site \(j\) has spin minus and the final spin is not plus, the carried
spin starts plus and cannot decrease under \(F(0)\). Its leading
term again vanishes, compensating for the possible simple pole in
the multiplier. This proves both refined bounds at every spectator.

\paragraph{Compatibility of cup reductions.}
To prove the adjacent reduction at \(x_{j+1}=Tx_j\) among the first
\(N-1\) sites, check it at the last-variable interpolation nodes.
Disjoint cup specializations commute by the slides and lower-size
exchange and reduction. Their scalar factors commute as well: for
two cup parameters \(a,d\), their cross factor is
\[
 \prod_{y\in\{d,Td\}}(y-t^5a)(y-t^6a)
 =T^2\prod_{r\in\{t^2,T,t^5,t^6\}}(d-ra).
\]
The four roots occur in reciprocal pairs and have product one, so
this expression is unchanged on exchanging \(a,d\).

At the overlapping node \(x_N=Tx_{j+1}\), both sides vanish: one
has the \(x_N\) spectator factor, and the other the \(x_j\) factor
in \(V_{j+1}\), whose slot is not crossed. At
\(x_N=Tx_j=x_{j+1}\), the first crossing in \(V_j\) is a checked
identity. It leaves the required cup and the lower-size exchange
transport of the remaining spectator. At the zero node in spin zero,
the scalar factors agree because \(t^{5+6}=T\). Interpolation proves
the adjacent reduction. Together with the exchanges already proved,
this gives every specialization \(x_k=Tx_i\), \(i<k\), with the
stated transport for nonadjacent slots.

\paragraph{The last exchange.}
Clear denominators in the remaining exchange discrepancy:
\[
 \mathcal D=(t^2x_{N-1}-x_N)(Tx_{N-1}-x_N)
 \big[F_{N-1,N}(x_{N-1}/x_N)\Psi(X)-\Psi(s_{N-1}X)\big].
\]
It vanishes on every cup hyperplane \(x_k=Tx_i\), \(i<k\), except
possibly the pair \((N-1,N)\). For an unaffected pair this is lower-size
exchange after cup reduction. When the second member is exchanged,
the cup transports acquire precisely that exchange or its inverse.
Consequently \(\mathcal D=Q E\), where
\[
 Q=\prod_{\substack{i<k\\(i,k)\ne(N-1,N)}}(x_k-Tx_i).
\]
Each spectator \(x_j\), \(j\le N-2\), occurs in exactly \(N-1\)
factors of \(Q\), exhausting the degree bound on \(\mathcal D\).
Thus \(E\) is independent of every spectator. The plus-spin degree
loss and minus-spin zero force all those spectator spins to be zero.
For adjacent spectators, however,
\(Q(s_jX)/Q(X)=b(x_j/x_{j+1})\). Their exchange equation therefore
forces \(F E=b E\), so their pair lies in channel 1. A nonzero pure
\(|00\rangle\) does not belong to \(\operatorname{im}I\), and there
are at least two spectators. Hence \(E=0\), proving full exchange.

\paragraph{Cyclicity.}
Lower-size cyclicity and the cup reductions show that
\[
 \Pi_\zeta^{-1}\Psi(\rho X)-\Psi(X)
 =\prod_{i<k\le N-1}(x_k-Tx_i)\,E(X).
\]
The specialized pair never straddles the rotation seam. The quotient
has degree at most one in each of its first \(N-1\) variables, so write
\[
 E(X)=\sum_{\epsilon\in\{0,1\}^{N-1}}
 x_1^{\epsilon_1}\cdots x_{N-1}^{\epsilon_{N-1}}
 E_\epsilon(x_N).
\]
By the refined bounds, bit zero permits spins \(0,+\) and bit one
permits spins \(0,-\). Exchange of adjacent variables among these
first \(N-1\) acts on \(E\) by
\[
 E(s_jX)=\big(aP_2+P_1+(a/b)P_0\big)_{j,j+1}E(X),
 \qquad a=a(x_j/x_{j+1}),\quad b=b(x_j/x_{j+1}).
\]

\begin{claim}
A tensor polynomial multiaffine in at least three consecutive variables,
with the bit-dependent spin support and exchange relation just stated,
is zero.
\end{claim}

\begin{proof}[Proof of the claim]
For an adjacent pair, the two uncanceled denominator factors of
\(a/b\), namely \(t^2x_j-x_{j+1}\) and \(Tx_{j+1}-x_j\),
must divide \(P_0E\). Their product has degree two in each
variable, whereas \(P_0E\) is multiaffine. Hence \(P_0E=0\).
The channel-1 part is symmetric. The channel-2 part is divisible by
\(x_{j+1}-t^2x_j\), with quotient independent of those two variables.
Consequently, with all other bits fixed, the adjacent coefficients obey
\[
 E_{10}=(P_1-t^2P_2)E_{01}.
\]
This operator is invertible on channels 1 and 2.

It suffices, therefore, to kill coefficients with sorted bits. Equal
adjacent bits require channel 1. With the permitted spins, their pair
must be proportional to \(I_+\) for 00 or \(I_-\) for 11.
Three equal bits are impossible: a first pair proportional to \(I_h\)
and either third spin \(0,h\) would produce a forbidden second pair
\(00\) or \(hh\). For 0011 the four slots must factor as
\(I_+\otimes I_-\). Their middle cup contraction is zero because
\(P_0E=0\), but fixing the end spins \(+,-\) makes that contraction
nonzero unless the coefficient vanishes. For 001, fix the first spin
plus and then zero. The cup contraction on the last two excludes
third spin zero and then minus, killing the coefficient through the
first-pair combination. The argument for 011 is the same with the
ends reversed. These exhaust sorted binary strings of length at least
three: either some bit occurs three times, both occur at least twice,
or the string is 001 or 011. This proves the claim.
\end{proof}

Apply the claim to the first \(N-1\) variables. Since \(N\ge4\),
the cyclic discrepancy vanishes.

\paragraph{Verification of the initial cases.}
The size-two and size-three formulas above meet the cup, zero, and
degree conditions directly. To check the first exchange by the explicit
blocks of \(F\), write \(S\) for the corresponding empty polynomial
and substitute \(e^{2iv}=a/b\) in the local weights. The scalar
identities used are
\[
\begin{aligned}
 P(b,a;c)&=w_1 P(a,b;c)+u_2 S/(2\sin\lambda\sqrt{ab}),\\
 P(c,b;a)&=u_1\sqrt{a/b}\,P(c,a;b)+o P(b,c;a),\\
 P(a,c;b)&=oP(c,a;b)+u_1\sqrt{b/a}\,P(b,c;a).
\end{aligned}
\]
(with square-root branch corresponding to \(v\)); at size 2 only the first, with the two \(P\)'s replaced by 1. For the first identities compare coefficients, using
\[
\begin{aligned}
 1-w_1&=u_2\cos v/\sin\lambda,\\
 \cos(2\lambda-v)-w_1\cos(2\lambda+v)
 &=(u_2/(2\sin\lambda))(\cos(2v)+\sqrt2).
\end{aligned}
\]
For the other two use \(u_1 e^{iv}+o/t=u_1 e^{-iv}+ot=\sin(3\lambda-v)/\sin(3\lambda+v)=(b+t a)/(a+t b)\). The \(w_2\) terms on the occupied pair cancel since \(t^2=i\); likewise the two contributions to the change of the empty component. Remaining exchanges follow by rotation not wrapping the pair. This verifies the bases, and finishes the construction.

\end{proof}

\subsubsection*{The physical normalization and its empty component}

The polynomial vector now determines the physical common eigenline.
We identify its empty component, which will be the denominator of the
separator pairing.

First,
\[
 \mathcal T_\zeta(z;X)\Psi(X)=\Psi(X).
\]
At \(z=x_N\) the last crossing is a checked identity. The row is
\(\Pi_\zeta\) followed by transport of its first parameter to the
end, so the exchange and cyclic equations prove the equality there.
Exchange gives it at all \(x_j\), row fusion/inversion gives it at
all \(Tx_j\), and triangularity gives it at zero. After clearing row
denominators, the degree is at most \(2N\). The \(2N+1\) generic
points \(0,x_1,\ldots,x_N,Tx_1,\ldots,Tx_N\) therefore suffice.

Similarly \(\Psi(X^{-1})^t\), with \(X^{-1}\) the entrywise reciprocal
list, is a left eigenvector of eigenvalue 1 for twist \(\zeta^{-1}\).
At \(z=x_N\), reverse transport through the symmetric crossings has
exactly the exchange arguments that move the last reciprocal parameter
to the first; the final rotation cancels because
\(\Pi_{\zeta^{-1}}^t\Pi_\zeta=1\). The other interpolation points
follow as above by inverse and transpose exchanges.

Let \(S_N\) be the empty component, which is not identically zero by
cup reduction, and put \(\Omega=\Psi/S_N\). For generic sites near
homogeneity, Lemma~\ref{lem:pv:physical-eigenline} identifies \(\Omega\)
with the lower half-cylinder state obtained from empty input far below,
normalized to have empty component 1. Its reciprocal transpose is the
analogously normalized left state at twist \(\zeta^{-1}\).

The symmetry of \(S_N\) requires this geometric normalization: it does
not follow just by taking the empty component of exchange, since that
component mixes with occupied pairs. Take real nearby angles with
\(\alpha_i>\alpha_{i+1}\) in the earlier embedding. From left to
right the upper boundary has site steps of directions
\(8\alpha_{i+1}/3,8\alpha_i/3\). Adding the rhombus above those
steps exchanges their order, entering from right and bottom with
parameters \(x_i,x_{i+1}\), without crossing the twist seam.
With empty final output, all gauge factors telescope and all loops
cancel even at finite height. The added crossing therefore takes one
normalized half-cylinder state to the other with no scalar factor.
Comparing with polynomial exchange gives \(S_N(s_iX)=S_N(X)\).
This extends polynomially from the open set of real angles.

\begin{proposition}[Empty component]\label{prop:pv:empty-component}
Put \(D(a,b)=a^2+\sqrt2 ab+b^2\) and
\(K(a,b)=(a^2-b^2)/D(a,b)\). The symmetric empty component is
\[
 S_N(X)=\prod_{i<j}\frac{D(x_i,x_j)}{x_i-x_j}\ {\rm Pf}[K(x_i,x_j)]
\]
where for odd size an extra last Pfaffian index has
\(K_{i,N+1}=1\), \(i\le N\), corresponding to parameter zero.
It satisfies
\[
 S_N(X^{-1})=\frac{S_N(X)}{\prod_i x_i^{N-1}}
\]
and the second fusion identity
\[
 S_N(a,t^2a,Y)=\sqrt2\,t a\prod_{y\in Y}(y-t^5 a)\ S_{N-1}(ta,Y).
\]
\end{proposition}

\begin{proof}
The displayed Pfaffian expression is a symmetric polynomial of the
required degrees. Its cleared numerator is alternating, so the
Vandermonde denominators cancel; the remaining possible poles cancel
against the factors \(D\). The bounded kernel at infinity gives the
one-site degree bound. At a first pair \((a,Ta)\), only its own
Pfaffian pairing survives and gives \(s_C S_{N-2}\). A last zero
parameter gives the zero-node prescription; if a dummy zero is present,
only their mutual pairing remains. Symmetry and interpolation identify
the expression with \(S_N\). Inversion of all parameters gives the
reciprocity formula. The zero-node identity and reciprocity also give
the leading-coefficient identity
\[
 [x_N^{N-1}]S_N(X)=S_{N-1}(x_1,\ldots,x_{N-1}).
\]

For the second fusion identity, compare its two sides in a spectator
variable. Cup reduction at another spectator is compatible on both
sides because
\[
 (a-t^5d)(a-t^6d)(t^2a-t^5d)(t^2a-t^6d)
 =(d-t^5a)(Td-t^5a)(ta-t^5d)(ta-t^6d)
\]
for a spectator pair \(d,Td\). Each of the other \(N-3\) spectators
therefore gives two roots, at ratios \(T\) and \(T^{-1}\), by symmetry.
The zero-node identity gives one more root. The leading-coefficient
identity makes the top degrees agree, lowering the degree of the
difference to at most \(N-2\). Thus its \(2(N-3)+1\) roots suffice
for \(N\ge4\); sizes two and three follow from the initial formulas.
\end{proof}

\begin{corollary}[Normalized reductions]\label{cor:pv:normalized-reductions}
At generic parameters, the normalized states reduce by insertion of
\(C\) at adjacent sites \((a,Ta)\), and by insertion of \(I\) at
\((a,t^2a)\), replacing those sites by the single site \(ta\).
When one parameter tends to zero or infinity, the spin-zero component
at that slot tends to the reduced \(\Omega\), with scalar factor one.
\end{corollary}

\begin{proof}
The cup assertion follows from the identical scalar reductions of
\(\Psi\) and \(S_N\). At \((a,t^2a)\), the pole in the exchange
equation forces channel 1. On that subspace the row intertwines the
reduced row by Lemma~\ref{lem:pv:row-relations}. Its eigenvalue-1
space is a unique line for generic reduced sites: it is one-dimensional
at homogeneous sites, and rational matrix minors preserve that property
generically. The reduced site \(ta\) can be generic. Empty-component
normalization fixes the insertion scalar to be one.

At zero, use rotation and the zero-node prescription. At infinity,
the limit exists generically by the degree bounds and has no plus
spin at the extreme slot. After projection to spin zero, triangularity
at crossing argument zero makes the transfer the reduced row.
Uniqueness again identifies the limit, and the nonzero generic leading
coefficient of \(S_N\) fixes its normalization.
\end{proof}
